\chapter{Esercizi Svolti}
\label{ch:esercizi}

Questo capitolo raccoglie esercizi svolti con soluzioni dettagliate per ciascuno degli argomenti trattati nella parte teorica. Si consiglia di provare a risolvere ogni esercizio autonomamente prima di consultare la soluzione.

%% =============================================================
\section{Esercizi su spazi vettoriali}
\label{sec:es-spazi}

\begin{examplebox}[Esercizio 1 --- Verifica di sottospazio]
Stabilire se $W = \{(x, y, z) \in \mathbb{R}^3 : x - 2y + z = 0\}$ è un sottospazio di $\mathbb{R}^3$.
\end{examplebox}

\textbf{Soluzione.} Verifichiamo le condizioni:
\begin{enumerate}
    \item \textit{Vettore nullo:} $(0, 0, 0)$ soddisfa $0 - 0 + 0 = 0$, quindi $\mathbf{0} \in W$.
    \item \textit{Chiusura:} Siano $\mathbf{u} = (u_1, u_2, u_3), \mathbf{v} = (v_1, v_2, v_3) \in W$ e $\lambda, \mu \in \mathbb{R}$. Allora $u_1 - 2u_2 + u_3 = 0$ e $v_1 - 2v_2 + v_3 = 0$. Per $\lambda\mathbf{u} + \mu\mathbf{v}$:
    \begin{align*}
    &(\lambda u_1{+}\mu v_1) - 2(\lambda u_2{+}\mu v_2) + (\lambda u_3{+}\mu v_3) \\
    &\quad = \lambda(u_1{-}2u_2{+}u_3) + \mu(v_1{-}2v_2{+}v_3) = 0.
    \end{align*}
\end{enumerate}
Quindi $W$ \textbf{è} un sottospazio di $\mathbb{R}^3$ (un piano per l'origine). \qed

\medskip

\begin{examplebox}[Esercizio 2 --- Non è un sottospazio]
Stabilire se $W = \{(x, y) \in \mathbb{R}^2 : xy \geq 0\}$ è un sottospazio di $\mathbb{R}^2$.
\end{examplebox}

\textbf{Soluzione.} $W$ contiene tutti i punti del primo e terzo quadrante. Consideriamo $\mathbf{u} = (1, 1) \in W$ e $\mathbf{v} = (-1, -1) \in W$ (entrambi hanno $xy > 0$). La somma:
\[ \mathbf{u} + \mathbf{v} = (1, 1) + (-1, -1) = (0, 0). \]
Questo va bene, ma prendiamo $\mathbf{u} = (1, 0) \in W$ e $\mathbf{v} = (0, -1) \in W$. La somma $\mathbf{u} + \mathbf{v} = (1, -1)$ ha $xy = -1 < 0$, quindi $\mathbf{u} + \mathbf{v} \notin W$.

$W$ \textbf{non} è un sottospazio (non è chiuso per la somma). \qed

\medskip

\begin{examplebox}[Esercizio 3 --- Base e dimensione]
Trovare una base e la dimensione del sottospazio $W = \text{Span}((1, 2, 3), (2, 4, 6), (1, 0, 1), (0, 1, 1))$ di $\mathbb{R}^3$.
\end{examplebox}

\textbf{Soluzione.} Costruiamo una matrice con i vettori come righe e riduciamo a scalini:
\[ \begin{pmatrix} 1 & 2 & 3 \\ 2 & 4 & 6 \\ 1 & 0 & 1 \\ 0 & 1 & 1 \end{pmatrix} \xrightarrow[\scriptscriptstyle R_3 - R_1]{\scriptscriptstyle R_2 - 2R_1} \begin{pmatrix} 1 & 2 & 3 \\ 0 & 0 & 0 \\ 0 & -2 & -2 \\ 0 & 1 & 1 \end{pmatrix} \xrightarrow[\scriptscriptstyle R_4 + \frac{1}{2}R_2]{\scriptscriptstyle R_2 \leftrightarrow R_3} \begin{pmatrix} 1 & 2 & 3 \\ 0 & -2 & -2 \\ 0 & 0 & 0 \\ 0 & 0 & 0 \end{pmatrix}. \]
Ci sono $2$ pivot, quindi $\dim W = 2$. Una base è data dai vettori corrispondenti ai pivot: $\{(1, 2, 3), (1, 0, 1)\}$ (oppure equivalentemente $\{(1, 2, 3), (0, 1, 1)\}$). \qed

\medskip

\begin{examplebox}[Esercizio 4 --- Intersezione e somma di sottospazi]
In $\mathbb{R}^3$, siano $W_1 = \text{Span}((1, 0, 1), (0, 1, -1))$ e $W_2 = \{(x,y,z) : x + y + z = 0\}$. Trovare $W_1 \cap W_2$, $W_1 + W_2$ e le loro dimensioni. Verificare la formula di Grassmann.
\end{examplebox}

\textbf{Soluzione.} Un generico elemento di $W_1$ è $\mathbf{v} = s(1,0,1) + t(0,1,-1) = (s, t, s-t)$. Imponendo $\mathbf{v} \in W_2$, cioè $s + t + (s - t) = 0$:
\[ 2s = 0 \quad \Longrightarrow \quad s = 0. \]
Quindi $W_1 \cap W_2 = \{(0, t, -t) : t \in \mathbb{R}\} = \text{Span}((0, 1, -1))$, con $\dim(W_1 \cap W_2) = 1$.

Per la somma: $\dim W_1 = 2$ (i generatori sono indipendenti), $\dim W_2 = 2$ (un'equazione in $\mathbb{R}^3$).

Formula di Grassmann: $\dim(W_1 + W_2) = 2 + 2 - 1 = 3 = \dim \mathbb{R}^3$. Quindi $W_1 + W_2 = \mathbb{R}^3$.

Verifica: $(1,0,1) \in W_1$ e $(0,1,-1) \in W_1 \cap W_2$. Un vettore di $W_2$ non in $W_1$: $(1, -1, 0) \in W_2$ (verifica: $1-1+0=0$). Questi tre vettori sono linearmente indipendenti (il determinante della matrice $3\times 3$ è non nullo), confermando che $W_1 + W_2 = \mathbb{R}^3$. \qed

\medskip

\begin{examplebox}[Esercizio 5 --- Coordinate rispetto a una base]
In $\mathbb{R}^3$, sia $\mathcal{B} = \{(1,1,0), (1,0,1), (0,1,1)\}$. Trovare le coordinate di $\mathbf{v} = (3, 2, 1)$ rispetto a $\mathcal{B}$.
\end{examplebox}

\textbf{Soluzione.} Cerchiamo $\lambda_1, \lambda_2, \lambda_3$ tali che $\lambda_1(1,1,0) + \lambda_2(1,0,1) + \lambda_3(0,1,1) = (3,2,1)$:
\[
\begin{cases}
\lambda_1 + \lambda_2 = 3 \\
\lambda_1 + \lambda_3 = 2 \\
\lambda_2 + \lambda_3 = 1
\end{cases}
\]
Sommando tutte: $2(\lambda_1 + \lambda_2 + \lambda_3) = 6$, quindi $\lambda_1 + \lambda_2 + \lambda_3 = 3$. Sottraendo ciascuna equazione:
\begin{align*}
    \lambda_3 &= 3 - 3 = 0, \\
    \lambda_2 &= 3 - 2 = 1, \\
    \lambda_1 &= 3 - 1 = 2.
\end{align*}
Quindi $[\mathbf{v}]_{\mathcal{B}} = \begin{pmatrix} 2 \\ 1 \\ 0 \end{pmatrix}$.

Verifica: $2(1,1,0) + 1(1,0,1) + 0(0,1,1) = (2,2,0) + (1,0,1) = (3,2,1)$. \checkmark \qed

%% =============================================================
\section[Matrici e sistemi]{Esercizi su matrici e sistemi lineari}
\label{sec:es-matrici}

\begin{examplebox}[Esercizio 6 --- Prodotto di matrici e inversa]
Date $A = \begin{pmatrix} 1 & 2 \\ 3 & 5 \end{pmatrix}$ e $B = \begin{pmatrix} 5 & -2 \\ -3 & 1 \end{pmatrix}$, verificare che $B = A^{-1}$.
\end{examplebox}

\textbf{Soluzione.} Calcoliamo $AB$:
\begin{align*}
AB &= \begin{pmatrix} 1 & 2 \\ 3 & 5 \end{pmatrix}\begin{pmatrix} 5 & -2 \\ -3 & 1 \end{pmatrix} = \begin{pmatrix} 5{-}6 & -2{+}2 \\ 15{-}15 & -6{+}5 \end{pmatrix} = \begin{pmatrix} -1 & 0 \\ 0 & -1 \end{pmatrix} = -I.
\end{align*}
Dunque $AB = -I$, cioè $B \neq A^{-1}$. In realtà $B = -A^{-1}$, e l'inversa corretta è:
\[ A^{-1} = -B = \begin{pmatrix} -5 & 2 \\ 3 & -1 \end{pmatrix}. \]
Verifica con la formula: $\det(A) = 5 - 6 = -1$, quindi $A^{-1} = \frac{1}{-1}\begin{pmatrix} 5 & -2 \\ -3 & 1 \end{pmatrix} = \begin{pmatrix} -5 & 2 \\ 3 & -1 \end{pmatrix}$. \checkmark \qed

\medskip

\begin{examplebox}[Esercizio 7 --- Sistema con parametro]
Risolvere il sistema al variare di $k \in \mathbb{R}$:
\[
\begin{cases}
x + y + kz = 1 \\
x + ky + z = 1 \\
kx + y + z = 1
\end{cases}
\]
\end{examplebox}

\textbf{Soluzione.} La matrice completa è:
\[ (A|\mathbf{b}) = \left(\begin{array}{ccc|c} 1 & 1 & k & 1 \\ 1 & k & 1 & 1 \\ k & 1 & 1 & 1 \end{array}\right) \xrightarrow[\scriptscriptstyle R_3 - kR_1]{\scriptscriptstyle R_2 - R_1} \left(\begin{array}{ccc|c} 1 & 1 & k & 1 \\ 0 & k{-}1 & 1{-}k & 0 \\ 0 & 1{-}k & 1{-}k^2 & 1{-}k \end{array}\right). \]
Notiamo che $R_2$: $0, (k-1), -(k-1), 0$, ovvero $(k-1)(0, 1, -1, 0)$.

\textbf{Caso $k = 1$}: $R_2$ e $R_3$ diventano righe nulle. Il sistema si riduce a $x + y + z = 1$, con $\text{rk} = 1$ e $3 - 1 = 2$ variabili libere:
\[ (x, y, z) = (1 - s - t, \; s, \; t), \quad s, t \in \mathbb{R}. \]

\textbf{Caso $k \neq 1$}: Dividiamo $R_2$ per $(k-1)$ e $R_3$ ha il fattore $(1-k) = -(k-1)$:
\[ \left(\begin{array}{ccc|c} 1 & 1 & k & 1 \\ 0 & 1 & -1 & 0 \\ 0 & -(k{-}1) & 1{-}k^2 & 1{-}k \end{array}\right) \xrightarrow{\scriptscriptstyle R_3 + (k-1)R_2} \left(\begin{array}{ccc|c} 1 & 1 & k & 1 \\ 0 & 1 & -1 & 0 \\ 0 & 0 & 2{-}k{-}k^2 & 1{-}k \end{array}\right). \]
L'elemento $(3,3)$: $1-k^2-(k-1) = 1-k^2-k+1 = 2-k-k^2 = -(k^2+k-2) = -(k+2)(k-1)$.

\textbf{Sotto-caso $k = -2$}: l'elemento $(3,3) = 0$ e l'elemento $(3,4) = 1-(-2) = 3 \neq 0$. Dunque $\text{rk}(A) = 2 < 3 = \text{rk}(A|\mathbf{b})$: il sistema è \textbf{impossibile}.

\textbf{Sotto-caso $k \neq 1, k \neq -2$}: $\text{rk}(A) = \text{rk}(A|\mathbf{b}) = 3$, soluzione unica.

Dalla riga 3: $z = \frac{1-k}{-(k+2)(k-1)} = \frac{-(1-k)}{(k+2)(1-k)} \cdot \frac{-1}{-1} = \frac{1}{k+2}$.

Dalla riga 2: $y = z = \frac{1}{k+2}$.

Dalla riga 1: $x = 1 - y - kz = 1 - \frac{1}{k+2} - \frac{k}{k+2} = 1 - \frac{k+1}{k+2} = \frac{1}{k+2}$.

\textbf{Riepilogo:}
\begin{itemize}
    \item $k \neq 1, k \neq -2$: soluzione unica $(x,y,z) = \left(\frac{1}{k+2}, \frac{1}{k+2}, \frac{1}{k+2}\right)$.
    \item $k = 1$: $\infty^2$ soluzioni: $(1-s-t, s, t)$.
    \item $k = -2$: impossibile. \qed
\end{itemize}

\medskip

\begin{examplebox}[Esercizio 8 --- Rango e Gauss]
Calcolare il rango di $A$ al variare di $\alpha \in \mathbb{R}$:
\[ A = \begin{pmatrix} 1 & 2 & \alpha \\ 2 & \alpha & 4 \\ 3 & 6 & \alpha+2 \end{pmatrix}. \]
\end{examplebox}

\textbf{Soluzione.} Riduciamo a scalini:
\[ \xrightarrow[R_3 \to R_3 - 3R_1]{R_2 \to R_2 - 2R_1} \begin{pmatrix} 1 & 2 & \alpha \\ 0 & \alpha - 4 & 4 - 2\alpha \\ 0 & 0 & 2 - 2\alpha \end{pmatrix}. \]
Notiamo che $R_3$: l'elemento $(3,3) = 2 - 2\alpha$ e $R_2$: l'elemento $(2,2) = \alpha - 4$.

\textbf{Caso $\alpha \neq 4$ e $\alpha \neq 1$}: tre pivot, $\text{rk}(A) = 3$.

\textbf{Caso $\alpha = 1$}: $R_3 = (0, 0, 0)$, $R_2 = (0, -3, 2)$, due pivot: $\text{rk}(A) = 2$.

\textbf{Caso $\alpha = 4$}: $R_2 = (0, 0, -4)$, $R_3 = (0, 0, -6)$. Dopo $R_3 \to R_3 - \frac{3}{2}R_2$: $R_3 = (0,0,0)$. Due pivot: $\text{rk}(A) = 2$.

In sintesi: $\text{rk}(A) = 3$ se $\alpha \neq 1$ e $\alpha \neq 4$; $\text{rk}(A) = 2$ se $\alpha = 1$ o $\alpha = 4$. \qed

%% =============================================================
\section{Esercizi sui determinanti}
\label{sec:es-determinanti}

\begin{examplebox}[Esercizio 9 --- Determinante $3 \times 3$]
Calcolare il determinante:
\[ A = \begin{pmatrix} 2 & -1 & 3 \\ 0 & 4 & -2 \\ 1 & 0 & 5 \end{pmatrix}. \]
\end{examplebox}

\textbf{Soluzione.} Sviluppiamo lungo la prima colonna (ha uno zero):
\begin{align*}
    \det(A) &= 2 \cdot \det\begin{pmatrix} 4 & -2 \\ 0 & 5 \end{pmatrix} - 0 + 1 \cdot (-1)^{3+1}\det\begin{pmatrix} -1 & 3 \\ 4 & -2 \end{pmatrix} \\
    &= 2(20 - 0) + 1(2 - 12) \\
    &= 40 - 10 = 30. \qed
\end{align*}

\medskip

\begin{examplebox}[Esercizio 10 --- Determinante $4 \times 4$]
Calcolare:
\[ \det\begin{pmatrix} 1 & 0 & 2 & 0 \\ 3 & 1 & 0 & 1 \\ 0 & 0 & 1 & 0 \\ 2 & 0 & 0 & 3 \end{pmatrix}. \]
\end{examplebox}

\textbf{Soluzione.} Sviluppiamo lungo la terza riga (tre zeri):
\[ \det(A) = (-1)^{3+3} \cdot 1 \cdot \det\begin{pmatrix} 1 & 0 & 0 \\ 3 & 1 & 1 \\ 2 & 0 & 3 \end{pmatrix}. \]
Sviluppiamo quest'ultimo lungo la prima riga:
\[ = 1 \cdot \det\begin{pmatrix} 1 & 1 \\ 0 & 3 \end{pmatrix} = 1 \cdot (3 - 0) = 3. \]
Quindi $\det(A) = 3$. \qed

\medskip

\begin{examplebox}[Esercizio 11 --- Inversa con l'aggiunta]
Calcolare l'inversa di $A = \begin{pmatrix} 1 & 0 & 1 \\ 0 & 2 & 0 \\ 1 & 0 & 3 \end{pmatrix}$.
\end{examplebox}

\textbf{Soluzione.} Calcoliamo $\det(A)$ sviluppando lungo la seconda riga:
\[ \det(A) = 2 \cdot (-1)^{2+2} \det\begin{pmatrix} 1 & 1 \\ 1 & 3 \end{pmatrix} = 2(3 - 1) = 4. \]
Matrice dei cofattori (calcolando tutti i $C_{ij}$):
\begin{align*}
    C_{11} &= \det\!\begin{psmallmatrix} 2 & 0 \\ 0 & 3 \end{psmallmatrix} = 6, \quad
    C_{12} = -\det\!\begin{psmallmatrix} 0 & 0 \\ 1 & 3 \end{psmallmatrix} = 0, \quad
    C_{13} = \det\!\begin{psmallmatrix} 0 & 2 \\ 1 & 0 \end{psmallmatrix} = -2, \\
    C_{21} &= -\det\!\begin{psmallmatrix} 0 & 1 \\ 0 & 3 \end{psmallmatrix} = 0, \quad
    C_{22} = \det\!\begin{psmallmatrix} 1 & 1 \\ 1 & 3 \end{psmallmatrix} = 2, \quad
    C_{23} = -\det\!\begin{psmallmatrix} 1 & 0 \\ 1 & 0 \end{psmallmatrix} = 0, \\
    C_{31} &= \det\!\begin{psmallmatrix} 0 & 1 \\ 2 & 0 \end{psmallmatrix} = -2, \quad
    C_{32} = -\det\!\begin{psmallmatrix} 1 & 1 \\ 0 & 0 \end{psmallmatrix} = 0, \quad
    C_{33} = \det\!\begin{psmallmatrix} 1 & 0 \\ 0 & 2 \end{psmallmatrix} = 2.
\end{align*}
La matrice aggiunta è la trasposta della matrice dei cofattori:
\[ \text{adj}(A) = \begin{pmatrix} 6 & 0 & -2 \\ 0 & 2 & 0 \\ -2 & 0 & 2 \end{pmatrix}. \]
Quindi:
\[ A^{-1} = \frac{1}{4}\begin{pmatrix} 6 & 0 & -2 \\ 0 & 2 & 0 \\ -2 & 0 & 2 \end{pmatrix} = \begin{pmatrix} 3/2 & 0 & -1/2 \\ 0 & 1/2 & 0 \\ -1/2 & 0 & 1/2 \end{pmatrix}. \]
Verifica: $AA^{-1} = I_3$. \checkmark \qed

\medskip

\begin{examplebox}[Esercizio 12 --- Cramer]
Risolvere con la regola di Cramer:
\[ \begin{cases} x + y + z = 6 \\ 2x - y + z = 3 \\ x + y - z = 0 \end{cases} \]
\end{examplebox}

\textbf{Soluzione.} Matrice dei coefficienti e suo determinante:
\[ A = \begin{pmatrix} 1 & 1 & 1 \\ 2 & -1 & 1 \\ 1 & 1 & -1 \end{pmatrix}, \quad \det(A) = 0 + 3 + 3 = 6. \]
Poiché $\det(A) = 6 \neq 0$, la soluzione è unica:
\begin{align*}
    x &= \tfrac{1}{6}\det\begin{psmallmatrix} 6 & 1 & 1 \\ 3 & -1 & 1 \\ 0 & 1 & -1 \end{psmallmatrix} = \frac{6(1{-}1) - 1(-3) + 1(3)}{6} = \frac{6}{6} = 1, \\[4pt]
    y &= \tfrac{1}{6}\det\begin{psmallmatrix} 1 & 6 & 1 \\ 2 & 3 & 1 \\ 1 & 0 & -1 \end{psmallmatrix} = \frac{-3+18-3}{6} = \frac{12}{6} = 2, \\[4pt]
    z &= \tfrac{1}{6}\det\begin{psmallmatrix} 1 & 1 & 6 \\ 2 & -1 & 3 \\ 1 & 1 & 0 \end{psmallmatrix} = \frac{-3+3+18}{6} = \frac{18}{6} = 3.
\end{align*}
La soluzione è $(x, y, z) = (1, 2, 3)$. \qed

%% =============================================================
\section[Applicazioni lineari]{Esercizi sulle applicazioni lineari}
\label{sec:es-applicazioni}

\begin{examplebox}[Esercizio 13 --- Nucleo, immagine e matrice associata]
Sia $T \colon \mathbb{R}^3 \to \mathbb{R}^3$ definita da $T(x,y,z) = (x+y, \; x-z, \; y+z)$. Trovare la matrice associata (basi canoniche), il nucleo, l'immagine e stabilire se $T$ è un isomorfismo.
\end{examplebox}

\textbf{Soluzione.} La matrice associata rispetto alle basi canoniche:
\[ A = \begin{pmatrix} 1 & 1 & 0 \\ 1 & 0 & -1 \\ 0 & 1 & 1 \end{pmatrix}. \]
\textbf{Nucleo:} Risolviamo $A\mathbf{v} = \mathbf{0}$:
\[ \begin{pmatrix} 1 & 1 & 0 \\ 1 & 0 & -1 \\ 0 & 1 & 1 \end{pmatrix} \xrightarrow{\scriptscriptstyle R_2 - R_1} \begin{pmatrix} 1 & 1 & 0 \\ 0 & -1 & -1 \\ 0 & 1 & 1 \end{pmatrix} \xrightarrow{\scriptscriptstyle R_3 + R_2} \begin{pmatrix} 1 & 1 & 0 \\ 0 & -1 & -1 \\ 0 & 0 & 0 \end{pmatrix}. \]
Da $R_2$: $y = -z$. Da $R_1$: $x = -y = z$. Quindi $\ker(T) = \{(t, -t, t) : t \in \mathbb{R}\} = \text{Span}((1,-1,1))$.

$\dim\ker(T) = 1$, quindi $T$ \textbf{non} è iniettiva e non è un isomorfismo.

\textbf{Immagine:} Per il teorema della dimensione, $\dim\text{Im}(T) = 3 - 1 = 2$. L'immagine è generata dalle colonne di $A$ che corrispondono ai pivot, ad esempio le prime due colonne: $\text{Im}(T) = \text{Span}((1,1,0), (1,0,1))$. \qed

\medskip

\begin{examplebox}[Esercizio 14 --- Determinare un'applicazione lineare]
Trovare l'applicazione lineare $T \colon \mathbb{R}^2 \to \mathbb{R}^2$ tale che $T(1,1) = (3, 2)$ e $T(1,-1) = (1, 0)$.
\end{examplebox}

\textbf{Soluzione.} Poiché $\{(1,1), (1,-1)\}$ è una base di $\mathbb{R}^2$, $T$ è determinata univocamente. Esprimiamo $(x,y)$ nella base data:
\[ (x,y) = \frac{x+y}{2}(1,1) + \frac{x-y}{2}(1,-1). \]
Per linearità:
\begin{align*}
    T(x,y) &= \frac{x+y}{2}T(1,1) + \frac{x-y}{2}T(1,-1) \\
    &= \frac{x+y}{2}(3,2) + \frac{x-y}{2}(1,0) \\
    &= \left(\frac{3(x+y) + (x-y)}{2}, \; \frac{2(x+y)}{2}\right) \\
    &= (2x + y, \; x + y).
\end{align*}
Verifica: $T(1,1) = (3, 2)$ \checkmark, $T(1,-1) = (1, 0)$ \checkmark.

Matrice (basi canoniche): $A = \begin{pmatrix} 2 & 1 \\ 1 & 1 \end{pmatrix}$. \qed

\medskip

\begin{examplebox}[Esercizio 15 --- Cambio di base]
Sia $T \colon \mathbb{R}^2 \to \mathbb{R}^2$ con matrice $A = \begin{pmatrix} 3 & 1 \\ 0 & 2 \end{pmatrix}$ rispetto alla base canonica $\mathcal{E}$. Trovare la matrice di $T$ rispetto alla base $\mathcal{B} = \{(1,1), (1,0)\}$.
\end{examplebox}

\textbf{Soluzione.} La matrice di cambio di base da $\mathcal{B}$ a $\mathcal{E}$ è:
\[ P = \begin{pmatrix} 1 & 1 \\ 1 & 0 \end{pmatrix} \quad \Longrightarrow \quad P^{-1} = \begin{pmatrix} 0 & 1 \\ 1 & -1 \end{pmatrix}. \]
La matrice di $T$ rispetto a $\mathcal{B}$ è:
\begin{align*}
    M_{\mathcal{B}}(T) &= P^{-1}AP = \begin{pmatrix} 0 & 1 \\ 1 & -1 \end{pmatrix}\begin{pmatrix} 3 & 1 \\ 0 & 2 \end{pmatrix}\begin{pmatrix} 1 & 1 \\ 1 & 0 \end{pmatrix} \\
    &= \begin{pmatrix} 0 & 2 \\ 3 & -1 \end{pmatrix}\begin{pmatrix} 1 & 1 \\ 1 & 0 \end{pmatrix} = \begin{pmatrix} 2 & 0 \\ 2 & 3 \end{pmatrix}. \qed
\end{align*}

%% =============================================================
\section[Autovalori]{Esercizi su autovalori e diagonalizzazione}
\label{sec:es-autovalori}

\begin{examplebox}[Esercizio 16 --- Autovalori e diagonalizzazione $3 \times 3$]
Diagonalizzare, se possibile, la matrice:
\[ A = \begin{pmatrix} 1 & 0 & 0 \\ 0 & 2 & 1 \\ 0 & 0 & 2 \end{pmatrix}. \]
\end{examplebox}

\textbf{Soluzione.} Polinomio caratteristico:
\[ p_A(\lambda) = \det\begin{pmatrix} 1-\lambda & 0 & 0 \\ 0 & 2-\lambda & 1 \\ 0 & 0 & 2-\lambda \end{pmatrix} = (1-\lambda)(2-\lambda)^2. \]
Autovalori: $\lambda_1 = 1$ con $m_a = 1$, e $\lambda_2 = 2$ con $m_a = 2$.

\textbf{Autospazio di $\lambda_1 = 1$:}
\[ A - I = \begin{pmatrix} 0 & 0 & 0 \\ 0 & 1 & 1 \\ 0 & 0 & 1 \end{pmatrix} \xrightarrow{R_2 \to R_2 - R_3} \begin{pmatrix} 0 & 0 & 0 \\ 0 & 1 & 0 \\ 0 & 0 & 1 \end{pmatrix}. \]
$y = 0$, $z = 0$, $x$ libero: $V_1 = \text{Span}((1,0,0))$, $m_g(1) = 1$. \checkmark

\textbf{Autospazio di $\lambda_2 = 2$:}
\[ A - 2I = \begin{pmatrix} -1 & 0 & 0 \\ 0 & 0 & 1 \\ 0 & 0 & 0 \end{pmatrix}. \]
$x = 0$, $z = 0$, $y$ libero: $V_2 = \text{Span}((0,1,0))$, $m_g(2) = 1 < 2 = m_a(2)$.

Poiché $m_g(2) \neq m_a(2)$, la matrice \textbf{non è diagonalizzabile}. \qed

\medskip

\begin{examplebox}[Esercizio 17 --- Diagonalizzazione completa]
Diagonalizzare:
\[ A = \begin{pmatrix} 1 & -1 & 0 \\ -1 & 2 & -1 \\ 0 & -1 & 1 \end{pmatrix}. \]
\end{examplebox}

\textbf{Soluzione.} La matrice è simmetrica, quindi è sicuramente diagonalizzabile (teorema spettrale).

Polinomio caratteristico:
\begin{align*}
    p_A(\lambda) &= \det\begin{pmatrix} 1-\lambda & -1 & 0 \\ -1 & 2-\lambda & -1 \\ 0 & -1 & 1-\lambda \end{pmatrix} \\
    &= (1-\lambda)\big[(2-\lambda)(1-\lambda) - 1\big] + 1\big[-(1-\lambda)\big] \\
    &= (1-\lambda)(2 - 3\lambda + \lambda^2 - 1) - (1-\lambda) \\
    &= (1-\lambda)(\lambda^2 - 3\lambda + 1) - (1-\lambda) \\
    &= (1-\lambda)(\lambda^2 - 3\lambda + 1 - 1) \\
    &= (1-\lambda) \cdot \lambda(\lambda - 3) \\
    &= -\lambda(\lambda - 1)(\lambda - 3).
\end{align*}
Autovalori: $\lambda_1 = 0$, $\lambda_2 = 1$, $\lambda_3 = 3$ (tutti distinti $\Rightarrow$ diagonalizzabile).

\textbf{Autospazio $V_0$} ($A\mathbf{v} = \mathbf{0}$):
\[ \begin{pmatrix} 1 & -1 & 0 \\ -1 & 2 & -1 \\ 0 & -1 & 1 \end{pmatrix} \to \begin{pmatrix} 1 & -1 & 0 \\ 0 & 1 & -1 \\ 0 & 0 & 0 \end{pmatrix}. \]
$y = z$, $x = y = z$: $V_0 = \text{Span}((1,1,1))$.

\textbf{Autospazio $V_1$} ($(A-I)\mathbf{v} = \mathbf{0}$):
\[ \begin{pmatrix} 0 & -1 & 0 \\ -1 & 1 & -1 \\ 0 & -1 & 0 \end{pmatrix} \to \begin{pmatrix} 1 & -1 & 1 \\ 0 & 1 & 0 \\ 0 & 0 & 0 \end{pmatrix}. \]
$y = 0$, $x = -z$: $V_1 = \text{Span}((1,0,-1))$.

\textbf{Autospazio $V_3$} ($(A-3I)\mathbf{v} = \mathbf{0}$):
\[ \begin{pmatrix} -2 & -1 & 0 \\ -1 & -1 & -1 \\ 0 & -1 & -2 \end{pmatrix} \to \begin{pmatrix} 1 & 0 & -1 \\ 0 & 1 & 2 \\ 0 & 0 & 0 \end{pmatrix}. \]
$x = z$, $y = -2z$: $V_3 = \text{Span}((1,-2,1))$.

La matrice diagonalizzante e la diagonale:
\[ P = \begin{pmatrix} 1 & 1 & 1 \\ 1 & 0 & -2 \\ 1 & -1 & 1 \end{pmatrix}, \quad P^{-1}AP = \begin{pmatrix} 0 & 0 & 0 \\ 0 & 1 & 0 \\ 0 & 0 & 3 \end{pmatrix}. \]
Poiché $A$ è simmetrica, si possono ortonormalizzare gli autovettori con Gram--Schmidt per ottenere una diagonalizzazione ortogonale. \qed

\medskip

\begin{examplebox}[Esercizio 18 --- Studio completo di diagonalizzabilità]
Studiare la diagonalizzabilità di $A$ al variare di $k \in \mathbb{R}$:
\[ A = \begin{pmatrix} 1 & k \\ 0 & 1 \end{pmatrix}. \]
\end{examplebox}

\textbf{Soluzione.} Polinomio caratteristico:
\[ p_A(\lambda) = (1 - \lambda)^2. \]
Unico autovalore: $\lambda = 1$ con $m_a = 2$.

Autospazio: $(A - I)\mathbf{v} = \mathbf{0}$:
\[ \begin{pmatrix} 0 & k \\ 0 & 0 \end{pmatrix}\begin{pmatrix} v_1 \\ v_2 \end{pmatrix} = \begin{pmatrix} 0 \\ 0 \end{pmatrix}. \]

\textbf{Se $k \neq 0$}: $v_2 = 0$, $v_1$ libero. $V_1 = \text{Span}((1,0))$, $m_g = 1 < 2 = m_a$. \textbf{Non diagonalizzabile}.

\textbf{Se $k = 0$}: $A = I$, che è già diagonale. $V_1 = \mathbb{R}^2$, $m_g = 2 = m_a$. \textbf{Diagonalizzabile} (anzi, $A$ è già diagonale). \qed

\medskip

\begin{examplebox}[Esercizio 19 --- Potenza di una matrice tramite diagonalizzazione]
Calcolare $A^{10}$ dove $A = \begin{pmatrix} 2 & 1 \\ 0 & 3 \end{pmatrix}$.
\end{examplebox}

\textbf{Soluzione.} $A$ è triangolare, quindi gli autovalori sono gli elementi diagonali: $\lambda_1 = 2$, $\lambda_2 = 3$ (distinti $\Rightarrow$ diagonalizzabile).

\textbf{Autovettori:}
\begin{itemize}
    \item $\lambda_1 = 2$: $(A - 2I)\mathbf{v} = \mathbf{0}$ dà $\begin{pmatrix} 0 & 1 \\ 0 & 1 \end{pmatrix}\mathbf{v} = \mathbf{0}$, cioè $v_2 = 0$. $\mathbf{v}_1 = (1,0)$.
    \item $\lambda_2 = 3$: $(A - 3I)\mathbf{v} = \mathbf{0}$ dà $\begin{pmatrix} -1 & 1 \\ 0 & 0 \end{pmatrix}\mathbf{v} = \mathbf{0}$, cioè $v_1 = v_2$. $\mathbf{v}_2 = (1,1)$.
\end{itemize}
Diagonalizzazione: $A = PDP^{-1}$ con
\[ P = \begin{pmatrix} 1 & 1 \\ 0 & 1 \end{pmatrix}, \quad D = \begin{pmatrix} 2 & 0 \\ 0 & 3 \end{pmatrix}, \quad P^{-1} = \begin{pmatrix} 1 & -1 \\ 0 & 1 \end{pmatrix}. \]
Quindi $A^{10} = PD^{10}P^{-1}$:
\begin{align*}
    A^{10} &= \begin{pmatrix} 1 & 1 \\ 0 & 1 \end{pmatrix}\begin{pmatrix} 2^{10} & 0 \\ 0 & 3^{10} \end{pmatrix}\begin{pmatrix} 1 & -1 \\ 0 & 1 \end{pmatrix} = \begin{pmatrix} 1024 & 3^{10} \\ 0 & 3^{10} \end{pmatrix}\begin{pmatrix} 1 & -1 \\ 0 & 1 \end{pmatrix} \\
    &= \begin{pmatrix} 1024 & 3^{10} - 1024 \\ 0 & 3^{10} \end{pmatrix} = \begin{pmatrix} 1024 & 58025 \\ 0 & 59049 \end{pmatrix}.
\end{align*}
Dove $3^{10} = 59049$. \qed

\medskip

\begin{examplebox}[Esercizio 20 --- Diagonalizzazione ortogonale]
Diagonalizzare ortogonalmente la matrice simmetrica:
\[ A = \begin{pmatrix} 3 & 1 \\ 1 & 3 \end{pmatrix}. \]
\end{examplebox}

\textbf{Soluzione.} Polinomio caratteristico: $p_A(\lambda) = (3-\lambda)^2 - 1 = \lambda^2 - 6\lambda + 8 = (\lambda-2)(\lambda-4)$.

Autovalori: $\lambda_1 = 2$, $\lambda_2 = 4$.

\textbf{Autovettori:}
\begin{itemize}
    \item $\lambda_1 = 2$: $(A-2I)\mathbf{v} = \mathbf{0}$ $\Rightarrow$ $\begin{pmatrix} 1 & 1 \\ 1 & 1 \end{pmatrix}\mathbf{v} = \mathbf{0}$ $\Rightarrow$ $v_1 = -v_2$. $\mathbf{v}_1 = (1,-1)$.
    \item $\lambda_2 = 4$: $(A-4I)\mathbf{v} = \mathbf{0}$ $\Rightarrow$ $\begin{pmatrix} -1 & 1 \\ 1 & -1 \end{pmatrix}\mathbf{v} = \mathbf{0}$ $\Rightarrow$ $v_1 = v_2$. $\mathbf{v}_2 = (1,1)$.
\end{itemize}
Verifica ortogonalità: $\mathbf{v}_1 \cdot \mathbf{v}_2 = 1 - 1 = 0$ \checkmark{} (come garantito dal teorema spettrale).

Normalizzazione: $\|\mathbf{v}_1\| = \|\mathbf{v}_2\| = \sqrt{2}$.
\[ Q = \frac{1}{\sqrt{2}}\begin{pmatrix} 1 & 1 \\ -1 & 1 \end{pmatrix}, \quad Q^TAQ = \begin{pmatrix} 2 & 0 \\ 0 & 4 \end{pmatrix}. \]
Verifica: $Q^TQ = I$ \checkmark{} (matrice ortogonale). \qed
